\subsection{Arithmetic certificates for a bounded packet}
The inverse compressions studied by Suzuki also give a complete finite
test at a prescribed height. Here are the spaces and the exact test.
Use $\widehat f(z)=\int_{\mathbb R}f(x)e^{-izx}\,dx$ and let
$\Gamma$ be the set of distinct signed ordinates of the zeta zeros.
Its multiplicity at $\gamma$ is $m_\gamma$. The Weil form is
\[
 q(f,g)=\sum_{\gamma\in\Gamma}
 m_\gamma\widehat f(\gamma)\overline{\widehat g(\gamma)}.
\]
Sampling identifies its completion with $\ell^2(\Gamma,m)$.
Let $C$ multiply the $\gamma$ coordinate by $1/\gamma$.
For an integer $N\ge1$, let $V_N$ consist of the continuous
piecewise-linear functions supported on $[-N,N]$, with knots
at multiples of $2^{-N}$, value zero at the endpoints, and integral
zero. Let $P_N$ be its $q$-orthogonal projection and $C_N=P_NCP_N$.
Define, for $k\ge1$,
\[
 S_k=-2k[z^{2k}]\log\frac{\Xi(z)}{\Xi(0)},\quad
 D_{N,k}=S_k-\operatorname{tr}C_N^{2k},\quad
 H_{N,k}=\left(\frac2{D_{N,k}}\right)^{1/(2k)}.
\]
Theorem~\ref{thm:fd-arithmetic-complete} proves that these defects
are positive, decrease in $N$, and satisfy
\[
 D_{N,k}\downarrow2\sum_{\gamma>0}(m_\gamma-1)\gamma^{-2k},
 \qquad
 \sum_{0<\gamma\le T}(m_\gamma-1)\le\tfrac12T^{2k}D_{N,k}.
\]
With $\gamma_*=x_*/2$, its complete bounded-height assertion is
\[
 \begin{split}
 &\text{every zero with }0<\Im\rho\le T\text{ is simple}\\
 &\quad\Longleftrightarrow
   \text{some finite }N,k\text{ satisfies }D_{N,k}<2/T^{2k},
 \end{split}
 \qquad
 \sup_{N,k}H_{N,k}=\gamma_*.
\]
In particular the Jensen limit is exactly
$\tfrac12(\sup_{N,k}H_{N,k})^2$.
The matrix entries are finite prime sums and archimedean integrals,
and the scalar $S_k$ is a Taylor coefficient of $\xi$ at $1/2$.
The interval calculation in Section~\ref{sec:fd-arithmetic} gives
\[
       \text{every zero with }0<\Im\rho\le97.632
       \text{ is simple}.
\]
Proposition~\ref{prop:fd-arithmetic-separation} also gives
$|\gamma_1-\gamma_2|>6\cdot10^{-10}$ for every pair of positive
zero occurrences through height $90$.
Platt--Trudgian established simplicity through height
$3\cdot10^{12}$ by rigorous sign counting and Turing's method
\cite[abstract, p.~792; preprint Section~2, pp.~2--3]{platt-trudgian2021}.
The examples here give the arithmetic matrices and divided-mass error
budgets for the present finite criteria.
The global assertion still asks that these certificates reach every
height. For the first trace defect, the orbital estimate gives
\[
 \lim_{a\downarrow0}\sqrt a\,m_*(-a)
       \ge\frac{\pi}{\sqrt{D_{N,1}}},\qquad
 \lim_{a\downarrow0}a\,j_*(-a)\ge\frac1{D_{N,1}}.
\]

There is also an exact integer certificate. Put
\(\iota_N(T)=\operatorname{ind}_+(T H_N-G_N)\), where
\(G_N\) and \(H_N\) are the Gram and compact-primitive matrices
of \(V_N\). At a height with \(\Xi(T)\ne0\),
Proposition~\ref{prop:fd-inertia-complete} proves
\[
 \iota_N(T)\uparrow n_{\rm dist}(T),
 \qquad
 \text{the packet through }T\text{ is simple}
 \ \Longleftrightarrow\ \exists N:\ \iota_N(T)\ge n(T).
\]
Here \(n\) counts multiplicity and \(n_{\rm dist}\) counts
support points. For every \(0<\eta<d_{\rm MT}\), where
\(d_{\rm MT}=5/4-\cot(1/\sqrt2)/(2\sqrt2)=0.8362503\ldots\),
Theorem~\ref{thm:strong-inertia-density} gives constants
\(A_\eta,T_\eta>0\) such that
\[
 \iota_{\lceil A_\eta\log T\rceil}(T)
       \ge(d_{\rm MT}-\eta)n(T)\qquad(T\ge T_\eta).
\]
Its matrices have dimension
\(O_\eta(T^{A_\eta\log2}\log T)\) and prime range
\(n<e^2T^{2A_\eta}\). The exact relation to the trace test is
\(D_{N,k}=4k\int_0^\infty(n(t)-\iota_N(t))t^{-2k-1}dt\).
The required zero integer deficit at unbounded heights remains open.

The elementary multiplicity bound $m_\gamma<(3/4)\log\gamma$
for $\gamma\ge10$, combined with Platt--Trudgian's certified
height $3\cdot10^{12}$, gives
\[
 \lim_{a\downarrow0}\sqrt a\,m_*(-a)
 =\lim_{a\downarrow0}\sqrt a\,R_*(-a)>1.65\cdot10^{12},
 \qquad \lim_{a\downarrow0}a\,j_*(-a)>4.5\cdot10^{24}.
\]
Corollary~\ref{cor:fd-global-growth-numeric} also gives an explicit
two-sided comparison of these global constants. The arithmetic trace
example through $97.632$ gives $216.88$ and $4766.003712$ by the
same maps, using its own finite matrix data.
Section~\ref{sec:executed-probe} carries out the spectral projection
and integer-rounding certificate through height $20$.
The contour certificate in Section~\ref{sec:executed-contour}
independently proves the same packet simple and reaches eight simple
centers through height $45$, with a common real heat interval
$|t|\le10^{-257}$. These statements use their separately certified
error budgets.

\subsection{Directions and common continuation of Hurwitz zeros}
For a primitive Dirichlet character $\chi$ of conductor $q$, including
the conductor-one case, use the denominator translation
\[
 F_\chi(u,s)=\sum_{n\ge1}\chi(n)(n+u)^{-s}
   =q^{-s}\sum_{r=1}^q\chi(r)\zeta\left(s,\frac{r+u}{q}\right).
\]
Continuation from $u=0$ specifies the branches. At a simple zero
$\rho=1/2+i\gamma$ of $L_\chi$, its tangent is
\[
                  v_\rho=\frac{\rho L_\chi(\rho+1)}{L_\chi'(\rho)}.
\]
For zeta, let $\theta$ be the continuous Riemann--Siegel phase with
$\theta(0)=0$, set $Z(t)=e^{i\theta(t)}\zeta(1/2+it)$, and define
$S$ by the midpoint convention $N(t)=\theta(t)/\pi+1+S(t)$.
At the $n$th positive simple ordinate put
\[
 W_n=S(\gamma_n)-\pi^{-1}\arg\zeta(3/2+i\gamma_n),\qquad
 \delta_n=\arctan\frac1{2\gamma_n},
\]
where the argument is the Euler-product logarithm.
Theorem~\ref{thm:hurwitz-sign} proves the exact identity
\[
 \Re v_{\rho_n}
 =\frac{|\zeta(\rho_n+1)|}{|Z'(\gamma_n)|}
      \sqrt{\gamma_n^2+1/4}\sin(\pi W_n+\delta_n).
\]
The rightward sign is therefore exactly the positivity of this sine.
At $n=1169$ the uncorrected sine is negative and the corrected sine
positive; both signs have rigorous interval certificates.
The velocity formula is due to Frisk--de Gosson
\cite[Section~2, equation~(3), p.~2]{frisk-degosson}.
The phase identity here fixes the midpoint and exterior-argument
conventions needed to use it.

Let $N_{\chi,T}$ count the zeros in $T<\Im\rho\le2T$ with
multiplicity. The signed-density theorem proves that, for each fixed
primitive $\chi$, each strict sign of $\Re v_\rho$ has lower
proportion at least $10^{-14}$ relative to $N_{\chi,T}$.
These counts concern distinct simple centers.
There are also absolute $\kappa,c,c_1,T_0>0$ such that, for
$T\ge T_0$ and $q\le T^\kappa$, at least
$10^{-30}N_{\chi,T}$ centers of each sign have simple, pairwise
disjoint branches with ordinates in $(T,2T)$ throughout
\[
 (1+T|u|)e^{12T|\Im u|}\le c_1e^{s_T},
                  \qquad s_T=\sqrt{\tfrac12\log\log T}.
\]
The sections are holomorphic in the interior and continue locally
across the boundary; their strict sides persist for real
$0<u\le c/T$. This is Theorem~\ref{thm:uniform-deformation}.
Theorems~\ref{thm:gaussian-speed-component} and
\ref{thm:anisotropic-continuation} give separate lower counts for
logarithmic derivative and velocity sizes, and larger continuation
regions obtained by selecting larger derivatives.

For the heat family, the same arithmetic coordinate instead gives
\[
 x'(0)=\frac{H_0''(x)}{H_0'(x)},\qquad
 s'(0)=\frac{iH_0''(x)}{2H_0'(x)}.
\]
These formulas specify the exact comparison of the two tangents.
At multiple Hurwitz centers and at simple vertical tangencies,
Section~\ref{sec:motion-bending} gives the complete second jet,
including the first nonzero normal coefficient of every multiple
branch.

\subsection{Backward heat survival and the isolated-pair law}
Put $M_T=\#\{x\in(T,2T):H_0(x)=0\}$, counting multiplicity.
Put $c_{\rm MT}=3/2-\cot(1/\sqrt2)/\sqrt2=0.6725007\ldots$,
the Montgomery--Taylor simple-zero proportion.
For every $0<\eta<c_{\rm MT}$ there are $c_\eta,T_\eta>0$ such that,
for $T\ge T_\eta$, at least $(c_{\rm MT}-\eta)M_T$ distinct simple
centers have branches that remain distinct, real, simple and in
$(T,2T)$ for
\[
                    -c_\eta/(\log T)^2\le t\le0.
\]
This is Theorem~\ref{thm:two-thirds-survival}. Its arithmetic input
is the restricted pair-correlation theorem, applied after the
zero-location theorem. Corollary~\ref{cor:fd-density-detection} proves that every
selected packet has infinite Toeplitz order, window size and Jensen
detection degree throughout that common interval. Multiplication by
this packet gives explicit comparison maps from the complementary
factor to the full heat function.
At the same time, absolute $C,a_0>0$ give, for $0<a<a_0$, a simple
nonreal zero of $H_{-a}$ with
\[
 0<\Re z\le\exp(\exp(C/\sqrt a)),\qquad\Im z\le-1/2.
\]
Theorem~\ref{thm:nonreal-height} uses Dobner's uniform deformed
Dirichlet-series approximation and an explicit finite recurrence
argument.

The local collision statement applies to a real entire function $f$
of order at most one whose zeros are all real. For two adjacent
simple zeros $b\pm\delta/2$, set
\[
 G(z)=\frac{f(z)}{(z-b)^2-\delta^2/4},\quad
 L_b=\frac{G'(b)}{G(b)},\quad
 S_b=\sum_{f(\xi)=0,\,\xi\ne b\pm\delta/2}(b-\xi)^{-2},
 \quad \kappa_b=\delta^2S_b.
\]
For sufficiently small $\kappa_b$, the local continuation under
$e^{a\partial_z^2}f$ has a unique real double collision near
$(a,z)=(\delta^2/8,b-\delta^2L_b/4)$, at
\[
 a_c=\delta^2\left(\frac18+\frac{\kappa_b}{32}
                         +O(\kappa_b^2)\right),\qquad
 z_c=b-2a_cL_b+O(\delta\kappa_b^{3/2}).
\]
The constants and the neighborhood after scaling by $\delta$ are
absolute. Theorem~\ref{thm:isolated-pair} proves this statement and
the creation of a conjugate nonreal pair. The quantity $\kappa_b$
measures exactly how the other zeros enter the error.

\subsection{The circular unitary model}
Let $\theta_1<\cdots<\theta_N$ be the angles of a Haar unitary
matrix, with $\theta_{i+N}=\theta_i+2\pi$, and put
\[
 Z_N(z)=\prod_{i=1}^N2\sin\frac{z-\theta_i}{2},\quad
 Z_{N,-a}=e^{a\partial_z^2}Z_N,
 \quad \delta_N=\min_i(\theta_{i+1}-\theta_i).
\]
For $N\ge2$ its first loss of reality $A_N$ is almost surely
positive and finite. Theorem~\ref{thm:cue-collisions} proves
\[
 \frac{8A_N}{\delta_N^2}=1+O_{\mathbb P}(N^{-1/3}),\qquad
 \mathbb P\left\{\frac{8N^{8/3}A_N}{(72\pi)^{2/3}}>v\right\}
                 \longrightarrow e^{-v^{3/2}}\quad(v\ge0).
\]
The first pair starts at the smallest gap with probability tending
to one. The collision time and location process, under the same
time scaling, converges on bounded time intervals to a Poisson
process of intensity $(3/2)\sqrt v\,dv\,db/(2\pi)$.
All collisions there are double with probability tending to one,
and each creates a nonreal pair that persists for the rest of the
interval. The number of such pairs at time $v$ therefore tends to
$\operatorname{Poisson}(v^{3/2})$.
Ben Arous--Bourgade prove the required extreme-gap process
\cite[Theorem~1.1 and Corollary~1.2, pp.~2650--2651]{ben-arous-bourgade};
the heat calculation transfers their gaps to these collision times.

For the label statement fix $r\in[0,1)$ and $a\in\mathbb R$ and
choose an independent uniform index $I$. Define
\[
 \phi_{N,r}(x)=\Im\sum_l\log(1-re^{i(x-\theta_l)}),\qquad
 W_{N,i}=i-\frac{N\theta_i}{2\pi}-a-\frac{\phi_{N,r}(\theta_i)}\pi,
 \qquad \eta_{N,i}=\operatorname{sgn}\sin(\pi W_{N,i}),
\]
with the power-series logarithm, and let
$D_{j,N}=N(\theta_{I+j}-\theta_I)/(2\pi)$.
Theorem~\ref{thm:cue-labels} proves that for every fixed $j\ge1$,
$(W_{N,I}\bmod2,D_{j,N})$ tends to $(V,D_j)$, where $V$ is
uniform modulo two and independent of the $j$th positive point of
the reduced Palm sine process. Its kernel is
$k(x-y)-k(x)k(y)$, where $k(u)=\sin(\pi u)/(\pi u)$.
Thus the labels are asymptotically fair and their limiting
correlation is
$\mathbb E[1-2d_2(j-D_j)]$, with
$d_2(v)=\min_{m\in\mathbb Z}|v-2m|$.
The convergence is uniform on compact ranges of $r<1$.
The limiting correlation tends to zero as $j\to\infty$.
For adjacent labels the limiting agreement probability lies strictly
between
\[
 0.6665466018372107200719020708703
 \quad\text{and}\quad
 0.6665466018372107200719020748216.
\]
The two independent interval calculations are described in
Section~\ref{sec:cue-labels}.
Section~\ref{sec:arithmetic-collisions} states the exact gap and
isolation estimates that would transfer the collision limit to the
zeta divisor, and gives a real-zero model showing the limitation of
the proved fixed-gap bounds.

\subsection{The Weil ground-state programme}
Let $A_L$ be the positive operator of the closed Weil form on
$(-L,L)$ against the $L^2$ norm, with eigenvalues
$a_1(L)\le a_2(L)\le\cdots$ counted with multiplicity.
Write $a_e,a_o$ for its even and odd minima.
Proposition~\ref{prop:parity-comparison} proves that
$a_e<a_o$ gives a simple even ground state $e_L$ with nonzero
integral and a Fourier transform with simple real zeros. Along
every sequence $L\to\infty$ with
$\inf(1-a_e/a_o)>0$, its normalized transform tends locally
uniformly to
\[
 \frac{\widehat e_L(z)}{\widehat e_L(0)}
       \longrightarrow\prod_{\gamma>0,\,\gamma\in\Gamma}
                         (1-z^2/\gamma^2).
\]
Each distinct ordinate contributes one factor. Suzuki's
Theorem~3.3, equation~(4.1), and Remark~4.1
\cite[pp.~10--12]{suzuki} give the unshifted inverse-compression
convergence. The fixed-multiplicity recovery by a probe in
Section~\ref{sec:suzuki-probes} is a direct consequence of that
theorem: the limiting spectral mass is
$m_\gamma|\widehat\varphi(\gamma)|^2$, and division and rounding
recover $m_\gamma$. The explicit error includes the reciprocal
spectral separation and the probe's Fourier lower bound.

Let $K$ be the inverse Fourier transform of $\Xi$.
For every family of real even $L^2$-unit vectors $v_L$ in the
local form domains with $v_L\to K/\|K\|_2$, the exact implication
proved in Proposition~\ref{prop:liminf-rayleigh} is
\[
 \liminf_{L\to\infty}\frac{q(v_L)}{a_2(L)}=0
       \quad\Longrightarrow\quad
       \text{every nontrivial zeta zero is simple}.
\]
The even prolate trial of Connes--Consani--Moscovici is defined
explicitly in Section~\ref{sec:prolate-trial}; for $L=\log\lambda$
it satisfies, for every $\varepsilon>0$,
\[
 q(v_L)\ll_\varepsilon\lambda^{9+\varepsilon}e^{-4\pi\lambda^2}.
\]
An odd trial tending to $K'/\|K'\|_2$ also proves
$a_2(L)\ll_\varepsilon\lambda^{13+\varepsilon}e^{-4\pi\lambda^2}$.
Thus the remaining lower estimate has absolute polynomial exponent
$9+\alpha$, with $0<\alpha\le4$; the parameter $\alpha$ measures
the gain above the even trial's exponent $9$.
There is a weaker all-window lower estimate, proved in
Theorem~\ref{thm:cosine-lower-rate}: for some absolute \(C,L_0>0\),
\[
 a_1(L)\ge\exp\{-C e^{2L/d_{\rm MT}}L(1+\log L)\}
                 \qquad(L\ge L_0).
\]
It also bounds \(a_2(L)\) and the codimension-one minimum \(b_L\)
of Section~\ref{sec:theta-coercivity}. Its exponent
\(1/d_{\rm MT}=1.1958141\ldots\) is strictly between one and \(6/5\).
Substitution into the trial ratio leaves a growing positive term of
order \(e^{2L/d_{\rm MT}}L\log L\) in the upper-bound exponent.
Thus the lower estimate at the displayed prolate scale remains open.

The exact compact-operator example at the end of the spectral
argument shows why the commutator and compactness estimates alone
leave this quantitative step open.
